\documentclass{article}
\usepackage[UTF8]{ctex}
\usepackage{amsmath,amssymb,amsfonts}
\usepackage{enumitem}
\usepackage{geometry}
\geometry{a4paper, margin=1in}
% === 错题本排版优化（不参与题目标识） ===
\setlength{\parindent}{0pt}
\setlength{\parskip}{0.35em plus 0.10em minus 0.08em}
\setlength{\emergencystretch}{3em}
\setlist{itemsep=0.25em, topsep=0.40em, parsep=0pt}
\allowdisplaybreaks[2]
% === 错题本排版优化结束 ===
\begin{document}
\section*{武汉大学2025年数学分析真题}
\section*{131．武汉大学2025年数学分析试卷}

\subsection{题目1}
1．求极限 \(\lim _{n \rightarrow \infty} \cos \frac{\pi}{2^{2}} \cos \frac{\pi}{2^{3}} \cdots \cos \frac{\pi}{2^{n}}\) ．

\subsection{题目2}
2．已知 \(\sin (x y)+2 y^{2}=1\) ，求 \(\frac{\mathrm{d} y}{\mathrm{~d} x}\) ．

\subsection{题目3}
3．\(I_{n}=\int_{0}^{\frac{\pi}{4}}(\sin x)^{n} \cos x \mathrm{~d} x\) ，求 \(\sum_{n=0}^{\infty} I_{n}\) ．

\subsection{题目4}
1．已知 \(D\) 是由 \(x=0, y=\frac{\pi}{2}, x=y^{2}\) 所围的区域，求 \(\iint_{D} \frac{\sin y}{y} \mathrm{~d} x \mathrm{~d} y\) ．

\subsection{题目5}
2．\(\sum\) 是 \(y=\sqrt{x}(0 \leq x \leq 1)\) 绕 \(x\) 轴旋转形成的曲面，法向量方向与 \(x\) 轴正方向成钝角，求第二型曲面积分

\[
I=\iint_{\sum} 2(x+y) \mathrm{d} y \mathrm{~d} z+y z \mathrm{~d} x \mathrm{~d} y
\]

\subsection{题目6}
3．已知 \(\varphi(t)=\int_{0}^{+\infty} \frac{\ln (1+t x)}{x(1+x)} \mathrm{d} x\) ，求 \(\varphi(1)\) ．

\subsection{题目7}
1．已知 \(p, q>0\) ，且 \(\frac{1}{p}+\frac{1}{q}=1\) ，求证：对任意正数 \(x_{1}, x_{2}\) ，有 \(x_{1} x_{2} \leq \frac{1}{p} x_{1}^{p}+\frac{1}{q} x_{2}^{q}\) ．

\subsection{题目8}
2．\(f(x)\) 在 \((-\infty,+\infty)\) 上二阶可微，且 \(f(x)>0, M=\sup _{x \in \mathrm{R}}\left|f^{\prime \prime}(x)\right|<+\infty\) 。证明：

\[
\left(f^{\prime}(x)\right)^{2} \leq 2 M f(x), x \in(-\infty,+\infty) .
\]

\subsection{题目9}
3．已知 \(f(x)\) 在 \([a, b]\) 上连续，并且满足对于 \(i=0,1,2\) ，有 \(\int_{a}^{b} x^{i} f(x) \mathrm{d} x=0\) ．求证：存在三个不同的数 \(x_{1}, x_{2}, x_{3} \in(a, b)\) ，使得 \(f\left(x_{1}\right)=f\left(x_{2}\right)=f\left(x_{3}\right)=0\) ．

\subsection{题目10}
4．\(\alpha\) 是正数，判断函数项级数 \(\sum_{n=0}^{\infty} x^{\alpha} e^{-n x}\) 在 \((0,+\infty)\) 上是否一致收敛。

\subsection{题目11}
5．已知 \(f(x, y)\) 在 \(x^{2}+y^{2}<1\) 上三阶连续可导，并且 \(f(0,0)=0\) ；求证：存在二阶连续可微函数\(g_{1}(x, y), g_{2}(x, y)\) ，使得 \(f(x, y)=x g_{1}(x, y)+y g_{2}(x, y)\) 。

\subsection{题目12}
6．求 \((a, b)\) 的两个区域，使得 \(f_{a, b}(x, y)=a y^{2}+b x\) 在约束条件 \(x^{2}+y^{2}=1\) 下分别有两个和四个临界点．
\end{document}
